\section{Introduction}

Modern systems increasingly construct task-specific filesystem views around
ordinary tools. Agents branch, fork, checkpoint, hide side effects, and run
build or test commands over repository
workspaces~\cite{agentfs_repo,branchfs_repo,yolofs_repo}. Environment systems
validate repositories against generated Docker state, cache epochs, and verified
local artifacts~\cite{docker_overview,swe_factory_repo,menvagent_repo,swe_rebench_v2_repo},
while services consume projected configuration, secrets, certificates, or
fixtures as operational state changes~\cite{kubernetes_projected_volumes,kubernetes_configmaps,kubernetes_secrets}.

However, these views make pathname binding a correctness boundary. Binding a
name to the wrong lower object can validate the wrong revision, leak a write
into a parent workspace, feed stale cache input into a build, or expose the
wrong service configuration. In many of these cases, the needed behavior is not
a new file format, storage service, or custom file operation. The needed
behavior chooses which existing object a name denotes, or whether that object is
visible.

The root cause is a boundary mismatch in today's filesystem extension points.
Workload state changes the right pathname-to-object relation at lookup or
directory-enumeration time, but ordinary file data, writes, permissions,
page-cache behavior, persistence, and consistency should remain with the lower
filesystem. A mechanism that acts before lookup must materialize or update a
view, while a mechanism that owns lookup as a filesystem service usually owns
more filesystem behavior than this subproblem needs.

Existing mechanisms bracket this lookup-time selection problem. Namespace
mechanisms such as bind mounts, OverlayFS, projected volumes, and copied trees
construct or materialize a namespace before the application uses
it~\cite{linux_mount,linux_overlayfs,kubernetes_projected_volumes}.
eBPF LSM hooks can attach verified access-control policy to security
decisions~\cite{linux_bpf_lsm},
but they are not an object-selection boundary. FUSE, stackable filesystems,
custom filesystems, and metadata services can compute richer views, but they
cross into a broader filesystem boundary~\cite{linux_fuse,wrapfs,bento,deltafs_sc21,indexfs,tablefs}.
They may be the right mechanism when the oracle needs synthetic contents,
custom metadata, data-path mediation, or runtime orchestration, but they are
broader than lookup-time name-resolution policy.

We argue that the missing narrow systems boundary is policy at VFS name
resolution, between access-control hooks and filesystem-service ownership. We
call the recurring subproblem a state-dependent path view, a workload-state
change that chooses which existing lower object a pathname denotes or whether
that object is visible. VFS name resolution is the point where
pathname-to-object selection and directory visibility are decided. The insight
is to separate that policy from filesystem ownership by letting a bounded
policy choose path-view actions at lookup and readdir time, then returning to
ordinary VFS and lower-filesystem semantics.

Making that separation practical creates three technical challenges. First, the
policy must run at the lookup and directory-enumeration operations that decide
the visible object, not in a separate namespace-construction phase. Second, the
policy output must be bounded and kernel-validated so the lower filesystem keeps
permissions, writes, data, page cache, persistence, and consistency. Third, the
same correctness condition must be usable to compare this boundary with a
feature-equivalent FUSE policy service and with the broader ownership required
by custom or stackable filesystems.

We present \namei as a \code{sched_ext}-style extension point for the VFS:
eBPF chooses name-resolution policy, while the kernel keeps the VFS mechanisms
that own objects and execute filesystem semantics~\cite{linux_sched_ext}. One
eBPF decision function receives lookup and directory-enumeration events,
returns bounded pass, redirect, hide, or kernel-registered directory-target
selection actions, and relies on kernel validation before ordinary VFS
execution resumes.
The policy does not implement file operations, allocate dentries or inodes, run
a filesystem daemon, or materialize a namespace tree. Together, these choices
answer the three challenges. The hook is placed at name resolution, the action
space is bounded and validated, and the mechanism preserves the same lower
objects used by the FUSE rows and referenced in the
custom/stackable-filesystem boundary account.

The evaluation follows the same boundary. RQ1 tests whether source-derived KVM
workloads pass the same correctness conditions while preserving lower-filesystem
semantics. RQ2 compares the same policies with feature-equivalent FUSE
implementations. RQ3 evaluates whether verifier-bounded policy and kernel
validation provide a narrower, fail-closed ownership boundary than custom or
stackable filesystems. The two primary workload evaluations are an Agent
workspace lifecycle and an environment/cache transition. Service/config is
included only if a concrete source-system correctness condition depends on
lookup-time object selection, so it remains motivating and conditional scope
rather than a primary evaluation.

The paper's primary contribution is the design and Linux implementation of
\namei as one systems boundary: a \code{sched_ext}-style VFS extension point
that places bounded policy at lookup and directory enumeration while preserving
VFS and lower-filesystem ownership. The paper makes three concrete
contributions: (1) the boundary design and policy contract for lookup and
readdir (Section~\ref{sec:design}); (2) a modified-Linux prototype with one
eBPF decision function, the real \code{cgroup/namei_ext} attachment path, and
kernel validation for bounded actions (Section~\ref{sec:implementation}); and
(3) an evaluation plan organized around the three RQs that fills correctness
results derived from source systems, feature-equivalent FUSE overhead, and
custom/stackable ownership-boundary result slots without treating workload
characterization as separate novelty (Section~\ref{sec:evaluation}).
